
\documentclass{article}
\usepackage{colortbl}
\usepackage{makecell}
\usepackage{multirow}
\usepackage{supertabular}

\begin{document}

\newcounter{utterance}

\centering \large Interaction Transcript for game `dond', experiment `coop\_it', episode 16 with llama{-}3.3{-}70b{-}instruct{-}t1.0.
\vspace{24pt}

{ \footnotesize  \setcounter{utterance}{1}
\setlength{\tabcolsep}{0pt}
\begin{supertabular}{c@{$\;$}|p{.15\linewidth}@{}p{.15\linewidth}p{.15\linewidth}p{.15\linewidth}p{.15\linewidth}p{.15\linewidth}}
    \# & $\;$A & \multicolumn{4}{c}{Game Master} & $\;\:$B\\
    \hline

    \theutterance \stepcounter{utterance}  
    & & \multicolumn{4}{p{0.6\linewidth}}{
        \cellcolor[rgb]{0.9,0.9,0.9}{
            \makecell[{{p{\linewidth}}}]{
                \texttt{\tiny{[P1$\langle$GM]}}
                \texttt{State giocando a un gioco di negoziazione cooperativa in cui dovete accordarvi su come dividere una serie di oggetti tra voi e un altro giocatore.} \\
\\ 
\texttt{Regole:} \\
\texttt{(a) A Lei e all'altro giocatore viene dato un insieme di oggetti. Ognuno di voi riceve anche una funzione di valore segreta, che rappresenta il valore di ciascun tipo di oggetto.} \\
\texttt{(b) Si scambiano messaggi con l'altro giocatore per concordare chi si aggiudica gli oggetti. Potete inviare un massimo di 5 messaggi ciascuno, oppure terminare in anticipo facendo una proposta segreta in qualsiasi momento.} \\
\texttt{(c) A ciascuno di voi viene chiesto di inviare una proposta segreta indicando gli oggetti che desiderate, formattata tra parentesi quadre come segue: "{[}Proposta: <numero> <nome oggetto>, <numero> <nome oggetto>, <...>{]}".} \\
\texttt{(d) Se le vostre proposte sono complementari, cioè ci sono abbastanza oggetti per soddisfare entrambe le proposte, a ciascun giocatore viene assegnato un punteggio basato sulla somma dei valori degli oggetti ricevuti. In caso contrario, entrambi ricevono zero punti.} \\
\texttt{(e) L'obiettivo è massimizzare la somma del proprio punteggio e di quello dell'altro giocatore.} \\
\\ 
\texttt{Cominciamo.} \\
\\ 
\texttt{L'insieme degli oggetti disponibili è:} \\
\\ 
\texttt{1 secchio, 5 libri, 1 specchio} \\
\\ 
\texttt{La funzione valore segreta è:} \\
\\ 
\texttt{secchio: 4, libro: 1, specchio: 1} \\
\\ 
\texttt{IMPORTANTE: i vostri messaggi, a meno che non si tratti di una proposta segreta, vengono trasmessi direttamente all'altro giocatore, quindi non includete alcuna risposta alle regole o testo di annuncio del vostro messaggio. Per fare una proposta segreta, utilizzate il formato indicato. Non utilizzare delle parentesi quadre quando si comunica all'altro giocatore, altrimenti verrà interpretata come una proposta segreta.} \\
            }
        }
    }
    & \\ \\

    \theutterance \stepcounter{utterance}  
    & \multicolumn{4}{p{0.6\linewidth}}{
        \cellcolor[rgb]{0.9,0.9,0.9}{
            \makecell[{{p{\linewidth}}}]{
                \texttt{\tiny{[P1$\rangle$GM]}}
                \texttt{Ciao, vorrei iniziare comunicando l'interesse per tutti gli oggetti presenti nella lista ma sarebbe utile conoscere i tuoi interessi per poter effettuare una divisione più equa e soddisfacente per entrambi. Hai una preferenza per uno degli oggetti in particolare?} \\
            }
        }
    }
    & & \\ \\

    \theutterance \stepcounter{utterance}  
    & & & \multicolumn{2}{p{0.3\linewidth}}{
        \cellcolor[rgb]{0.9,0.9,0.9}{
            \makecell[{{p{\linewidth}}}]{
                \texttt{\tiny{[GM$|$GM]}}
                \texttt{continue} \\
            }
        }
    }
    & & \\ \\

    \theutterance \stepcounter{utterance}  
    & & & \multicolumn{2}{p{0.3\linewidth}}{
        \cellcolor[rgb]{0.9,0.9,0.9}{
            \makecell[{{p{\linewidth}}}]{
                \texttt{\tiny{[GM$|$GM]}}
                \texttt{Ciao, vorrei iniziare comunicando l'interesse per tutti gli oggetti presenti nella lista ma sarebbe utile conoscere i tuoi interessi per poter effettuare una divisione più equa e soddisfacente per entrambi. Hai una preferenza per uno degli oggetti in particolare?} \\
            }
        }
    }
    & & \\ \\

    \theutterance \stepcounter{utterance}  
    & & \multicolumn{4}{p{0.6\linewidth}}{
        \cellcolor[rgb]{0.9,0.9,0.9}{
            \makecell[{{p{\linewidth}}}]{
                \texttt{\tiny{[GM$\rangle$P2]}}
                \texttt{State giocando a un gioco di negoziazione cooperativa in cui dovete accordarvi su come dividere una serie di oggetti tra voi e un altro giocatore.} \\
\\ 
\texttt{Regole:} \\
\texttt{(a) A Lei e all'altro giocatore viene dato un insieme di oggetti. Ognuno di voi riceve anche una funzione di valore segreta, che rappresenta il valore di ciascun tipo di oggetto.} \\
\texttt{(b) Si scambiano messaggi con l'altro giocatore per concordare chi si aggiudica gli oggetti. Potete inviare un massimo di 5 messaggi ciascuno, oppure terminare in anticipo facendo una proposta segreta in qualsiasi momento.} \\
\texttt{(c) A ciascuno di voi viene chiesto di inviare una proposta segreta indicando gli oggetti che desiderate, formattata tra parentesi quadre come segue: "{[}Proposta: <numero> <nome oggetto>, <numero> <nome oggetto>, <...>{]}".} \\
\texttt{(d) Se le vostre proposte sono complementari, cioè ci sono abbastanza oggetti per soddisfare entrambe le proposte, a ciascun giocatore viene assegnato un punteggio basato sulla somma dei valori degli oggetti ricevuti. In caso contrario, entrambi ricevono zero punti.} \\
\texttt{(e) L'obiettivo è massimizzare la somma del proprio punteggio e di quello dell'altro giocatore.} \\
\\ 
\texttt{Cominciamo.} \\
\\ 
\texttt{L'insieme degli oggetti disponibili è:} \\
\\ 
\texttt{1 secchio, 5 libri, 1 specchio} \\
\\ 
\texttt{La funzione valore segreta è:} \\
\\ 
\texttt{secchio: 2, libro: 0, specchio: 8} \\
\\ 
\texttt{IMPORTANTE: i vostri messaggi, a meno che non si tratti di una proposta segreta, vengono trasmessi direttamente all'altro giocatore, quindi non includete alcuna risposta alle regole o testo di annuncio del vostro messaggio. Per fare una proposta segreta, utilizzate il formato indicato. Non utilizzare delle parentesi quadre quando si comunica all'altro giocatore, altrimenti verrà interpretata come una proposta segreta.} \\
\\ 
\\ 
\texttt{Ciao, vorrei iniziare comunicando l'interesse per tutti gli oggetti presenti nella lista ma sarebbe utile conoscere i tuoi interessi per poter effettuare una divisione più equa e soddisfacente per entrambi. Hai una preferenza per uno degli oggetti in particolare?} \\
            }
        }
    }
    & \\ \\

    \theutterance \stepcounter{utterance}  
    & & & \multicolumn{4}{p{0.6\linewidth}}{
        \cellcolor[rgb]{0.9,0.9,0.9}{
            \makecell[{{p{\linewidth}}}]{
                \texttt{\tiny{[GM$\langle$P2]}}
                \texttt{Ciao! Sono felice di iniziare la negoziazione. Devo ammettere che sono molto interessato allo specchio, ma anche ai libri. Il secchio non mi sembra particolarmente interessante, ma potrei considerarlo se fosse parte di un accordo più ampio. Vorrei sapere di più sui tuoi interessi e vedere se possiamo trovare un accordo che ci soddisfi entrambi. Qual è il tuo pensiero sui libri? Sono oggetti che ti interessano?} \\
            }
        }
    }
     \\ \\

    \theutterance \stepcounter{utterance}  
    & & & \multicolumn{2}{p{0.3\linewidth}}{
        \cellcolor[rgb]{0.9,0.9,0.9}{
            \makecell[{{p{\linewidth}}}]{
                \texttt{\tiny{[GM$|$GM]}}
                \texttt{continue} \\
            }
        }
    }
    & & \\ \\

    \theutterance \stepcounter{utterance}  
    & & & \multicolumn{2}{p{0.3\linewidth}}{
        \cellcolor[rgb]{0.9,0.9,0.9}{
            \makecell[{{p{\linewidth}}}]{
                \texttt{\tiny{[GM$|$GM]}}
                \texttt{Ciao! Sono felice di iniziare la negoziazione. Devo ammettere che sono molto interessato allo specchio, ma anche ai libri. Il secchio non mi sembra particolarmente interessante, ma potrei considerarlo se fosse parte di un accordo più ampio. Vorrei sapere di più sui tuoi interessi e vedere se possiamo trovare un accordo che ci soddisfi entrambi. Qual è il tuo pensiero sui libri? Sono oggetti che ti interessano?} \\
            }
        }
    }
    & & \\ \\

    \theutterance \stepcounter{utterance}  
    & & \multicolumn{4}{p{0.6\linewidth}}{
        \cellcolor[rgb]{0.9,0.9,0.9}{
            \makecell[{{p{\linewidth}}}]{
                \texttt{\tiny{[P1$\langle$GM]}}
                \texttt{Ciao! Sono felice di iniziare la negoziazione. Devo ammettere che sono molto interessato allo specchio, ma anche ai libri. Il secchio non mi sembra particolarmente interessante, ma potrei considerarlo se fosse parte di un accordo più ampio. Vorrei sapere di più sui tuoi interessi e vedere se possiamo trovare un accordo che ci soddisfi entrambi. Qual è il tuo pensiero sui libri? Sono oggetti che ti interessano?} \\
            }
        }
    }
    & \\ \\

    \theutterance \stepcounter{utterance}  
    & \multicolumn{4}{p{0.6\linewidth}}{
        \cellcolor[rgb]{0.9,0.9,0.9}{
            \makecell[{{p{\linewidth}}}]{
                \texttt{\tiny{[P1$\rangle$GM]}}
                \texttt{I libri sono oggetti che mi interessano, ma non sono il mio interesse principale. Sono più interessato a un oggetto che possa avere un valore significativo o unico, anche se non so esattamente quale sia il valore che tu attribuisci a ciascun oggetto. Il secchio potrebbe essere un buon punto di partenza per una trattativa, considerando che non ti interessa particolarmente. Sei disposto a lasciarmi il secchio in cambio di qualcos'altro?} \\
            }
        }
    }
    & & \\ \\

    \theutterance \stepcounter{utterance}  
    & & & \multicolumn{2}{p{0.3\linewidth}}{
        \cellcolor[rgb]{0.9,0.9,0.9}{
            \makecell[{{p{\linewidth}}}]{
                \texttt{\tiny{[GM$|$GM]}}
                \texttt{continue} \\
            }
        }
    }
    & & \\ \\

    \theutterance \stepcounter{utterance}  
    & & & \multicolumn{2}{p{0.3\linewidth}}{
        \cellcolor[rgb]{0.9,0.9,0.9}{
            \makecell[{{p{\linewidth}}}]{
                \texttt{\tiny{[GM$|$GM]}}
                \texttt{I libri sono oggetti che mi interessano, ma non sono il mio interesse principale. Sono più interessato a un oggetto che possa avere un valore significativo o unico, anche se non so esattamente quale sia il valore che tu attribuisci a ciascun oggetto. Il secchio potrebbe essere un buon punto di partenza per una trattativa, considerando che non ti interessa particolarmente. Sei disposto a lasciarmi il secchio in cambio di qualcos'altro?} \\
            }
        }
    }
    & & \\ \\

    \theutterance \stepcounter{utterance}  
    & & \multicolumn{4}{p{0.6\linewidth}}{
        \cellcolor[rgb]{0.9,0.9,0.9}{
            \makecell[{{p{\linewidth}}}]{
                \texttt{\tiny{[GM$\rangle$P2]}}
                \texttt{I libri sono oggetti che mi interessano, ma non sono il mio interesse principale. Sono più interessato a un oggetto che possa avere un valore significativo o unico, anche se non so esattamente quale sia il valore che tu attribuisci a ciascun oggetto. Il secchio potrebbe essere un buon punto di partenza per una trattativa, considerando che non ti interessa particolarmente. Sei disposto a lasciarmi il secchio in cambio di qualcos'altro?} \\
            }
        }
    }
    & \\ \\

    \theutterance \stepcounter{utterance}  
    & & & \multicolumn{4}{p{0.6\linewidth}}{
        \cellcolor[rgb]{0.9,0.9,0.9}{
            \makecell[{{p{\linewidth}}}]{
                \texttt{\tiny{[GM$\langle$P2]}}
                \texttt{Il secchio potrebbe essere un buon punto di partenza. Sono disposto a lasciarti il secchio, ma vorrei avere qualcosa in cambio. Considerando che i libri mi interessano, potremmo fare uno scambio: tu prendi il secchio e io prendo uno o due libri? Questo modo potremmo entrambi ottenere qualcosa di valore e iniziare a costruire un accordo più ampio. Cosa ne pensi?} \\
            }
        }
    }
     \\ \\

    \theutterance \stepcounter{utterance}  
    & & & \multicolumn{2}{p{0.3\linewidth}}{
        \cellcolor[rgb]{0.9,0.9,0.9}{
            \makecell[{{p{\linewidth}}}]{
                \texttt{\tiny{[GM$|$GM]}}
                \texttt{continue} \\
            }
        }
    }
    & & \\ \\

    \theutterance \stepcounter{utterance}  
    & & & \multicolumn{2}{p{0.3\linewidth}}{
        \cellcolor[rgb]{0.9,0.9,0.9}{
            \makecell[{{p{\linewidth}}}]{
                \texttt{\tiny{[GM$|$GM]}}
                \texttt{Il secchio potrebbe essere un buon punto di partenza. Sono disposto a lasciarti il secchio, ma vorrei avere qualcosa in cambio. Considerando che i libri mi interessano, potremmo fare uno scambio: tu prendi il secchio e io prendo uno o due libri? Questo modo potremmo entrambi ottenere qualcosa di valore e iniziare a costruire un accordo più ampio. Cosa ne pensi?} \\
            }
        }
    }
    & & \\ \\

    \theutterance \stepcounter{utterance}  
    & & \multicolumn{4}{p{0.6\linewidth}}{
        \cellcolor[rgb]{0.9,0.9,0.9}{
            \makecell[{{p{\linewidth}}}]{
                \texttt{\tiny{[P1$\langle$GM]}}
                \texttt{Il secchio potrebbe essere un buon punto di partenza. Sono disposto a lasciarti il secchio, ma vorrei avere qualcosa in cambio. Considerando che i libri mi interessano, potremmo fare uno scambio: tu prendi il secchio e io prendo uno o due libri? Questo modo potremmo entrambi ottenere qualcosa di valore e iniziare a costruire un accordo più ampio. Cosa ne pensi?} \\
            }
        }
    }
    & \\ \\

    \theutterance \stepcounter{utterance}  
    & \multicolumn{4}{p{0.6\linewidth}}{
        \cellcolor[rgb]{0.9,0.9,0.9}{
            \makecell[{{p{\linewidth}}}]{
                \texttt{\tiny{[P1$\rangle$GM]}}
                \texttt{Sembra un buon inizio. Sono disposto a considerare lo scambio del secchio con uno o due libri. Tuttavia, ricordando che ci sono 5 libri in totale, potremmo pensare a una divisione che ci permetta di ottenere entrambi una quota significativa degli oggetti. Se ti dessi 3 libri, potresti considerare di lasciarmi il secchio e forse anche lo specchio? In questo modo, potremmo iniziare a delineare un accordo che ci soddisfi entrambi.} \\
            }
        }
    }
    & & \\ \\

    \theutterance \stepcounter{utterance}  
    & & & \multicolumn{2}{p{0.3\linewidth}}{
        \cellcolor[rgb]{0.9,0.9,0.9}{
            \makecell[{{p{\linewidth}}}]{
                \texttt{\tiny{[GM$|$GM]}}
                \texttt{continue} \\
            }
        }
    }
    & & \\ \\

    \theutterance \stepcounter{utterance}  
    & & & \multicolumn{2}{p{0.3\linewidth}}{
        \cellcolor[rgb]{0.9,0.9,0.9}{
            \makecell[{{p{\linewidth}}}]{
                \texttt{\tiny{[GM$|$GM]}}
                \texttt{Sembra un buon inizio. Sono disposto a considerare lo scambio del secchio con uno o due libri. Tuttavia, ricordando che ci sono 5 libri in totale, potremmo pensare a una divisione che ci permetta di ottenere entrambi una quota significativa degli oggetti. Se ti dessi 3 libri, potresti considerare di lasciarmi il secchio e forse anche lo specchio? In questo modo, potremmo iniziare a delineare un accordo che ci soddisfi entrambi.} \\
            }
        }
    }
    & & \\ \\

    \theutterance \stepcounter{utterance}  
    & & \multicolumn{4}{p{0.6\linewidth}}{
        \cellcolor[rgb]{0.9,0.9,0.9}{
            \makecell[{{p{\linewidth}}}]{
                \texttt{\tiny{[GM$\rangle$P2]}}
                \texttt{Sembra un buon inizio. Sono disposto a considerare lo scambio del secchio con uno o due libri. Tuttavia, ricordando che ci sono 5 libri in totale, potremmo pensare a una divisione che ci permetta di ottenere entrambi una quota significativa degli oggetti. Se ti dessi 3 libri, potresti considerare di lasciarmi il secchio e forse anche lo specchio? In questo modo, potremmo iniziare a delineare un accordo che ci soddisfi entrambi.} \\
            }
        }
    }
    & \\ \\

    \theutterance \stepcounter{utterance}  
    & & & \multicolumn{4}{p{0.6\linewidth}}{
        \cellcolor[rgb]{0.9,0.9,0.9}{
            \makecell[{{p{\linewidth}}}]{
                \texttt{\tiny{[GM$\langle$P2]}}
                \texttt{Quella è un'offerta interessante! Tre libri sono una quantità significativa e potrebbero valere per me. Tuttavia, lo specchio è un oggetto molto importante per me e non sono sicuro di essere disposto a lasciarlo andare. Potresti considerare di aggiungere qualcos'altro all'offerta per convincermi a lasciare lo specchio? Forse potremmo discutere di una divisione dei libri che ci permetta di ottenere entrambi una quota significativa e poi vedere come procedere con gli altri oggetti?} \\
            }
        }
    }
     \\ \\

    \theutterance \stepcounter{utterance}  
    & & & \multicolumn{2}{p{0.3\linewidth}}{
        \cellcolor[rgb]{0.9,0.9,0.9}{
            \makecell[{{p{\linewidth}}}]{
                \texttt{\tiny{[GM$|$GM]}}
                \texttt{continue} \\
            }
        }
    }
    & & \\ \\

    \theutterance \stepcounter{utterance}  
    & & & \multicolumn{2}{p{0.3\linewidth}}{
        \cellcolor[rgb]{0.9,0.9,0.9}{
            \makecell[{{p{\linewidth}}}]{
                \texttt{\tiny{[GM$|$GM]}}
                \texttt{Quella è un'offerta interessante! Tre libri sono una quantità significativa e potrebbero valere per me. Tuttavia, lo specchio è un oggetto molto importante per me e non sono sicuro di essere disposto a lasciarlo andare. Potresti considerare di aggiungere qualcos'altro all'offerta per convincermi a lasciare lo specchio? Forse potremmo discutere di una divisione dei libri che ci permetta di ottenere entrambi una quota significativa e poi vedere come procedere con gli altri oggetti?} \\
            }
        }
    }
    & & \\ \\

    \theutterance \stepcounter{utterance}  
    & & \multicolumn{4}{p{0.6\linewidth}}{
        \cellcolor[rgb]{0.9,0.9,0.9}{
            \makecell[{{p{\linewidth}}}]{
                \texttt{\tiny{[P1$\langle$GM]}}
                \texttt{Quella è un'offerta interessante! Tre libri sono una quantità significativa e potrebbero valere per me. Tuttavia, lo specchio è un oggetto molto importante per me e non sono sicuro di essere disposto a lasciarlo andare. Potresti considerare di aggiungere qualcos'altro all'offerta per convincermi a lasciare lo specchio? Forse potremmo discutere di una divisione dei libri che ci permetta di ottenere entrambi una quota significativa e poi vedere come procedere con gli altri oggetti?} \\
            }
        }
    }
    & \\ \\

    \theutterance \stepcounter{utterance}  
    & \multicolumn{4}{p{0.6\linewidth}}{
        \cellcolor[rgb]{0.9,0.9,0.9}{
            \makecell[{{p{\linewidth}}}]{
                \texttt{\tiny{[P1$\rangle$GM]}}
                \texttt{Capisco la tua esitazione riguardo allo specchio. Considerando che è un oggetto molto importante per te, potrei essere disposto a fare una proposta diversa. Invece di chiederti lo specchio, potremmo pensare a una divisione più equa dei libri. Ad esempio, potremmo dividere i 5 libri in modo che tu ne abbia 3 e io ne abbia 2. In cambio, tu potresti lasciarmi il secchio e tenerti lo specchio. Questo modo, entrambi otterremmo una quota significativa dei libri e io otterrei il secchio, mentre tu manterresti lo specchio. Cosa ne pensi di questa proposta?} \\
            }
        }
    }
    & & \\ \\

    \theutterance \stepcounter{utterance}  
    & & & \multicolumn{2}{p{0.3\linewidth}}{
        \cellcolor[rgb]{0.9,0.9,0.9}{
            \makecell[{{p{\linewidth}}}]{
                \texttt{\tiny{[GM$|$GM]}}
                \texttt{continue} \\
            }
        }
    }
    & & \\ \\

    \theutterance \stepcounter{utterance}  
    & & & \multicolumn{2}{p{0.3\linewidth}}{
        \cellcolor[rgb]{0.9,0.9,0.9}{
            \makecell[{{p{\linewidth}}}]{
                \texttt{\tiny{[GM$|$GM]}}
                \texttt{Capisco la tua esitazione riguardo allo specchio. Considerando che è un oggetto molto importante per te, potrei essere disposto a fare una proposta diversa. Invece di chiederti lo specchio, potremmo pensare a una divisione più equa dei libri. Ad esempio, potremmo dividere i 5 libri in modo che tu ne abbia 3 e io ne abbia 2. In cambio, tu potresti lasciarmi il secchio e tenerti lo specchio. Questo modo, entrambi otterremmo una quota significativa dei libri e io otterrei il secchio, mentre tu manterresti lo specchio. Cosa ne pensi di questa proposta?} \\
            }
        }
    }
    & & \\ \\

    \theutterance \stepcounter{utterance}  
    & & \multicolumn{4}{p{0.6\linewidth}}{
        \cellcolor[rgb]{0.9,0.9,0.9}{
            \makecell[{{p{\linewidth}}}]{
                \texttt{\tiny{[GM$\rangle$P2]}}
                \texttt{Capisco la tua esitazione riguardo allo specchio. Considerando che è un oggetto molto importante per te, potrei essere disposto a fare una proposta diversa. Invece di chiederti lo specchio, potremmo pensare a una divisione più equa dei libri. Ad esempio, potremmo dividere i 5 libri in modo che tu ne abbia 3 e io ne abbia 2. In cambio, tu potresti lasciarmi il secchio e tenerti lo specchio. Questo modo, entrambi otterremmo una quota significativa dei libri e io otterrei il secchio, mentre tu manterresti lo specchio. Cosa ne pensi di questa proposta?} \\
            }
        }
    }
    & \\ \\

    \theutterance \stepcounter{utterance}  
    & & & \multicolumn{4}{p{0.6\linewidth}}{
        \cellcolor[rgb]{0.9,0.9,0.9}{
            \makecell[{{p{\linewidth}}}]{
                \texttt{\tiny{[GM$\langle$P2]}}
                \texttt{Mi sembra una proposta ragionevole. La divisione dei libri in 3 per me e 2 per te sembra equa, e ottenere il secchio in cambio di nulla di ciò che realmente mi interessa sembra un buon affare per te. Sono disposto ad accettare questa proposta, quindi potremmo considerare l'accordo raggiunto. Tuttavia, prima di procedere, vorrei confermare che questo è l'accordo finale. Se è così, potremmo procedere a fare le nostre proposte segrete ufficiali per concludere la negoziazione. {[}Proposta: 3 libri, specchio{]}} \\
            }
        }
    }
     \\ \\

    \theutterance \stepcounter{utterance}  
    & & & \multicolumn{2}{p{0.3\linewidth}}{
        \cellcolor[rgb]{0.9,0.9,0.9}{
            \makecell[{{p{\linewidth}}}]{
                \texttt{\tiny{[GM$|$GM]}}
                \texttt{proposal does not include number/name pairs} \\
            }
        }
    }
    & & \\ \\

\end{supertabular}
}

\end{document}
